\documentclass{article}
\usepackage[T1]{fontenc}
\usepackage{lmodern}
\usepackage{graphicx}
\usepackage{amsmath,amssymb}
\usepackage[a4paper,margin=2cm]{geometry}
\usepackage[absolute,overlay]{textpos}
\setlength{\TPHorizModule}{1mm}
\setlength{\TPVertModule}{1mm}
\usepackage[indonesian]{babel}
\usepackage{hyperref}
\setlength{\emergencystretch}{2em}
% unit: Halaman 46

\begin{document}

% === Halaman 46 (python/native) ===

DES Berganda

• Karena DES mempunyai potensi kelemahan pada brute force atack, maka 

dibuat varian dari DES.

• Varian DES yang paling luas digunakan adalah DES berganda (multiple DES).

• DES berganda adalah enkripsi berkali-kali dengan DES dan menggunakan 

kunci ganda.

• Tinjau dua macam DES berganda:

 1. Double DES

 2. Triple DES (3DES)

46

\end{document}
